\Needspace{18\baselineskip}
\section{Recursive isomorphism and substitution modulo renaming}
\label{sec:recursive-isomorphism}

Literal equality of edge labels is too restrictive when a label contains a
whole graph. Renaming the body of a box changes its graph record, although its
mathematical incidence structure is unchanged. This section defines recursive
isomorphism, proves its invariants, and proves congruence of replacement at any
addressed depth. All comparisons are within a fixed depth bound and fixed
ordered outer interface. The quotient construction below does not identify
different depth bounds or prove the representation theorem for the paper's
set-based presentation.

\subsection{Isomorphisms with related labels}

\begin{definition}[Isomorphism relative to a label relation]\label{term:relational-isomorphism}
Fix a signature and a relation $R$ on its labels. An $R$-isomorphism from $G$ to
$H$ consists of bijections $v:G_V\to H_V$ of internal vertices and
$e:G_E\to H_E$ of edges. Internal vertex labels agree, corresponding edge labels
are related by $R$, and mapping every source or target list through the induced
port map gives the corresponding list in $H$. The port map fixes each ordered
input and output position. We write $G\cong_R H$ for the existence of such data.
\end{definition}
\lean{Graph/RelIso}{10}{16}{Net.RelIso}

\begin{lemma}[Identity, inverse, and composition]
If $R$ is reflexive, every graph has an identity $R$-isomorphism. If $R$ is
symmetric, every $R$-isomorphism has an inverse. If $R$ is transitive,
$R$-isomorphisms compose.
\end{lemma}
\begin{proof}
For identity use identity carrier maps and reflexivity of $R$. For inverse use
the inverse carrier bijections. Evaluate the original label equation at the
inverse image, cancel the edge bijection, and use symmetry of $R$. For incidence,
map the original incidence equation through the inverse port map; cancellation
recovers the reversed equation. Composition uses composed carrier bijections,
transitivity of the label relation, and the fact that a list mapped twice is
mapped by the composite. Port maps compose because they act on internal names
and fix both kinds of boundary position.
\end{proof}
\lean{Graph/RelIso}{22}{28}{Net.RelIso.refl}
\lean{Graph/RelIso}{30}{44}{Net.RelIso.symm}
\lean{Graph/RelIso}{46}{62}{Net.RelIso.trans}

\begin{lemma}[Incidence paths and ordinary isomorphisms]
Every $R$-isomorphism preserves and reflects adjacency, hence preserves and
reflects paths under its edge bijection. Every ordinary label-preserving
isomorphism gives an $R$-isomorphism when $R$ is reflexive.
\end{lemma}
\begin{proof}
Map a shared vertex to obtain forward adjacency. For reflection, source and
target preservation express the shared target vertex as images of two original
ports. Injectivity of the port bijection makes those ports equal, giving the
original shared vertex. Apply the existing list-path mapping lemma to this
adjacency equivalence. For an ordinary isomorphism, retain all carrier and
incidence data; equality of its labels and reflexivity give the required
$R$-relation at each edge.
\end{proof}
\lean{Graph/RelIso}{64}{78}{Net.RelIso.adjacent_iff}
\lean{Graph/RelIso}{80}{82}{Net.RelIso.path_iff}
\lean{Graph/RelIso}{86}{93}{Net.RelIso.ofIso}

\subsection{Recursive comparison of box contents}

\begin{definition}[Related boxes]\label{term:related-boxes}
For a relation on bodies of each fixed interface, two boxes are related when
they have the same number of bound variables, the same number of free outputs,
and related bodies at the interface specified by those two numbers. The
constructor $\mathsf{BoxRelated.bodies}$ enforces both metadata equalities; it
does not merely require equality of their sum.
\end{definition}
\lean{Graph/RecursiveIso}{10}{13}{Linear.BoxRelated}

\begin{definition}[Recursive label equivalence]\label{term:recursive-label-equivalence}
At depth zero, label equivalence is literal equality. At positive depth,
primitive labels are equivalent exactly when they are equal; a primitive and a
box are never equivalent. Two box labels are equivalent when their metadata
agree and there exists an isomorphism of their bodies whose corresponding
labels are recursively equivalent at the smaller depth. This is a well-founded
definition because the depth decreases whenever a box body is compared.
\end{definition}
\lean{Graph/RecursiveIso}{17}{23}{Linear.LabelEquiv}

\begin{definition}[Recursive hypernet isomorphism]\label{term:deep-isomorphism}
A recursive isomorphism $G\cong_d H$, represented by $\mathsf{DeepIso}$, is an
isomorphism relative to recursive label equivalence at depth $d$. Its outer
carrier bijections are explicit data. The label proof at a box supplies the
existence of a further body isomorphism, independently of the names chosen at
any other box. Every level preserves its own ordered boundary positions.
\end{definition}
\lean{Graph/RecursiveIso}{27}{27}{Linear.DeepIso}

\begin{lemma}[Equivalence laws for recursive labels]
Recursive label equivalence is reflexive, symmetric, and transitive.
\end{lemma}
\begin{proof}
Induct on depth. At depth zero the claims are equality laws. At positive depth,
primitive cases again reduce to equality and mixed primitive/box cases are
impossible. For box reflexivity, the induction hypothesis supplies reflexivity
of the body-label relation, so the identity relative isomorphism witnesses body
equivalence. For box symmetry, take a witnessing body isomorphism and invert it
using symmetry at the smaller depth. For transitivity, matching the two
box-relation constructors identifies all metadata, so the two body isomorphisms
have a common middle graph and interface. Compose them using transitivity at
the smaller depth. Package each resulting body isomorphism in the box relation.
\end{proof}
\lean{Graph/RecursiveIso}{30}{39}{Linear.labelEquiv_refl}
\lean{Graph/RecursiveIso}{42}{57}{Linear.labelEquiv_symm}
\lean{Graph/RecursiveIso}{60}{82}{Linear.labelEquiv_trans}
Identity, inverse, and composition of recursive graph isomorphisms follow by
instantiating the generic constructions with these three laws. Ordinary
isomorphisms embed into recursive ones by reflexivity of the label relation.
\lean{Graph/RecursiveIso}{87}{87}{Linear.DeepIso.refl}
\lean{Graph/RecursiveIso}{89}{89}{Linear.DeepIso.symm}
\lean{Graph/RecursiveIso}{91}{92}{Linear.DeepIso.trans}
\lean{Graph/RecursiveIso}{95}{95}{Linear.DeepIso.ofIso}

\begin{lemma}[Signatures and operation labels are invariant]
Equivalent labels have equal ordered input signatures, equal ordered output
signatures, and the same operation-label predicate.
\end{lemma}
\begin{proof}
At depth zero or for two primitives, substitute the label equality. Mixed
labels are impossible. For two boxes, the relation constructor identifies the
bound/free metadata, and both signatures depend only on those numbers. Both
box labels fail the operation predicate. No information about body carrier
names is used in these conclusions.
\end{proof}
\lean{Graph/RecursiveIso}{100}{112}{Linear.labelEquiv_input}
\lean{Graph/RecursiveIso}{115}{127}{Linear.labelEquiv_output}
\lean{Graph/RecursiveIso}{130}{142}{Linear.labelEquiv_isOperation}

\Needspace{16\baselineskip}
\subsection{Operation traces and paths}

\begin{theorem}[Recursive renaming preserves and reflects operation paths]
Let $f:G\cong_d H$, with port map $v$ and edge bijection $e$. An operation trace
from $x$ to $y$ along the edge list $\ell$ in $G$ exists exactly when the trace
from $v(x)$ to $v(y)$ along $\mathsf{map}(e,\ell)$ exists in $H$. The analogous
statement holds for forward operation paths.
\end{theorem}
\begin{proof}
First prove forward trace transport by induction on the trace. An empty trace
maps to an empty trace. For a nonempty trace, operation-label invariance supplies
the required label predicate. Source and target preservation map its two
incidence witnesses, and induction supplies the remaining trace.

For reflection transport the target trace through the inverse recursive
isomorphism. The inverse port and edge maps cancel the original maps, recovering
both original endpoints and exactly the original edge list. A forward operation
path is the reverse-oriented trace with its edge list reversed. List mapping
commutes with reversal, so the trace equivalence gives the path equivalence.
In particular, input and output endpoints retain their boundary indices.
\end{proof}
\lean{Graph/RecursiveIso}{145}{156}{Linear.DeepIso.operationTrace}
\lean{Graph/RecursiveIso}{159}{171}{Linear.DeepIso.operationTrace_iff}
\lean{Graph/RecursiveIso}{174}{181}{Linear.DeepIso.operationPath_iff}

\subsection{Transporting occurrences and holes through all boxes}

\begin{definition}[Fibre and hole embeddings]\label{term:hole-embedding}
A fibre embedding maps positions in one label to positions in another,
injectively, and preserves their complete hole observations, including the
absence of a hole. A hole embedding similarly maps graph occurrences
injectively and preserves their observations. These maps concern occurrence
addresses; their definition does not itself assert any incidence property.
\end{definition}
\lean{Graph/RecursiveIsoHoles}{9}{12}{Linear.FiberEmbedding}
\lean{Graph/RecursiveIsoHoles}{15}{18}{Linear.HoleEmbedding}

\begin{lemma}[Assembling fibre embeddings]
An injective map of outer edges and a fibre embedding over each mapped edge
together give a hole embedding of graphs.
\end{lemma}
\begin{proof}
Map $(a,x)$ to the mapped edge paired with the image of $x$ under its fibre
embedding. If two results coincide, projecting their edges and using edge-map
injectivity identifies the original edges. Their remaining fibre positions now
have the same type, and fibre-map injectivity identifies them. Observation
preservation is exactly the preservation property of the corresponding fibre
embedding.
\end{proof}
\lean{Graph/RecursiveIsoHoles}{21}{32}{Linear.HoleEmbedding.ofFibers}

\termdef{all-edge-transport}
A hole embedding also transports the original deep addresses: first use the
occurrence-coordinate bijection of the source, apply the hole embedding, then
use the inverse coordinate bijection of the target. This defines
$\mathsf{HoleEmbedding.allEdge}$.
\lean{Graph/RecursiveIsoHoles}{35}{37}{Linear.HoleEmbedding.allEdge}
The resulting address map is injective. Apply the target coordinate bijection
to an equality of images, cancel its inverse, use embedding injectivity, and
finally use source-coordinate injectivity. It preserves hole lookup because
both coordinate bijections preserve observation and the middle embedding does
so as well.
\lean{Graph/RecursiveIsoHoles}{39}{45}{Linear.HoleEmbedding.allEdge_injective}
\lean{Graph/RecursiveIsoHoles}{47}{53}{Linear.HoleEmbedding.allEdge_hole}

\begin{lemma}[Recursive labels admit fibre embeddings]
If two labels are recursively equivalent, an injective, observation-preserving
map of their fibres exists.
\end{lemma}
\begin{proof}
Induct on depth. Equal primitives use the identity map. Mixed label cases are
impossible. For equivalent boxes, choose a witnessing isomorphism of their
bodies. For each body edge, the induction hypothesis supplies a fibre embedding
from its recursive label-equivalence proof; choose one such map per edge.
Assemble these maps over the body isomorphism's injective edge map to obtain a
body hole embedding, and express it in original deep-address coordinates.

Map the box's shallow position to the shallow position, and each child address
through this body map. The two alternatives remain disjoint. Within the shallow
alternative equality is immediate; within the child alternative injectivity is
the proved injectivity of body-address transport. Shallow observations are empty
on both sides, while child observations agree by the body's transport theorem.
This constructs the desired embedding. Classical choice selects witnesses of
the existential body-isomorphism and embedding statements; it does not supply
any graph invariant as an additional assumption.
\end{proof}
\lean{Graph/RecursiveIsoHoles}{57}{96}{Linear.labelEquiv_hasEmbedding}
The noncomputable function $\mathsf{labelEquiv\_embedding}$ chooses the embedding
whose existence was just proved. Applying it at every outer edge of a recursive
isomorphism, and assembling over its edge bijection, defines
$\mathsf{DeepIso.holeEmbedding}$.
\lean{Graph/RecursiveIsoHoles}{99}{101}{Linear.labelEquiv_embedding}
\lean{Graph/RecursiveIsoHoles}{104}{107}{Linear.DeepIso.holeEmbedding}

\begin{lemma}[Presence transports and uniqueness reflects]
A hole embedding transports every present hole label from source to target. If
the target has unique hole observations, so does the source.
\end{lemma}
\begin{proof}
Map the occurrence witnessing a present label and use preservation of lookup.
For uniqueness, map two source occurrences with equal hole observations. Target
uniqueness makes their images equal, and embedding injectivity makes the
original occurrences equal.
\end{proof}
\lean{Graph/RecursiveIsoHoles}{110}{113}{Linear.HoleEmbedding.hasHole}
\lean{Graph/RecursiveIsoHoles}{116}{122}{Linear.HoleEmbedding.unique}

\begin{theorem}[Hole and context invariance under recursive isomorphism]
Recursively isomorphic graphs have identical hole inventories. One has unique
holes if and only if the other does; one is hole-free if and only if the other
is.
\end{theorem}
\begin{proof}
Construct hole embeddings from the isomorphism and its inverse. Presence
transport gives the two implications for inventory equality. Uniqueness
reflection, in the two directions, gives the context equivalence. Negating hole
presence for every label gives the hole-freeness equivalence. The two chosen
embeddings are not asserted to be mutually inverse; the two implications and
injectivity are sufficient for these results.
\end{proof}
\lean{Graph/RecursiveIsoHoles}{125}{127}{Linear.DeepIso.hasHole_iff}
\lean{Graph/RecursiveIsoHoles}{130}{132}{Linear.DeepIso.uniqueHoles_iff}
\lean{Graph/RecursiveIsoHoles}{134}{137}{Linear.DeepIso.holeFree_iff}

\subsection{Congruence of shallow and deep replacement}

\termdef{relational-plugging}
For a fixed host $C$ and shallow selected edge $h$, an $R$-isomorphism
$f:H\cong_R K$ gives the candidate carrier bijections for the plugged graphs:
use identity on the host summand and $f$ on the replacement summand. For edges,
the host summand is the subtype of edges different from $h$.
\lean{Graph/RelIsoPlug}{16}{17}{Net.RelIsoPlug.vertices}
\lean{Graph/RelIsoPlug}{19}{20}{Net.RelIsoPlug.edges}
The induced port map commutes with host embedding, by a case split on internal,
input, and output ports. It commutes with insertion after the replacement's own
port map: an internal port uses the right summand, and a boundary port reduces
to the host-embedding equation at the selected source or target port.
\lean{Graph/RelIsoPlug}{22}{24}{Net.RelIsoPlug.map_host}
\lean{Graph/RelIsoPlug}{26}{32}{Net.RelIsoPlug.map_insert}

\begin{theorem}[Shallow congruence for related labels]
For reflexive $R$, an $R$-isomorphism $H\cong_R K$ induces
$C[h:=H]\cong_R C[h:=K]$. No non-self-loop condition is required.
\end{theorem}
\begin{proof}
Use the preceding sum bijections. Host internal labels are unchanged and
replacement internal labels are preserved by $f$. Host edge labels satisfy
reflexivity of $R$, and inserted edge labels use the relation supplied by $f$.
For a host incidence list, commute the port map with host embedding pointwise.
For an inserted incidence list, commute it with insertion and then use the
source or target preservation equation of $f$. This establishes every field of
the resulting relative isomorphism without requiring insertion to be globally
injective.
\end{proof}
\lean{Graph/RelIsoPlug}{38}{85}{Net.RelIso.plug}

\begin{lemma}[Box congruence]
Boxing recursively isomorphic bodies with the same metadata gives recursively
isomorphic boxes. Replacing one fixed box edge in a surrounding graph by either
of those bodies also gives recursively isomorphic graphs.
\end{lemma}
\begin{proof}
For a single boxed graph, use identity maps on its empty internal carrier and
one-edge carrier. Its labels are related by the body-isomorphism witness, and
its boundary lists depend only on the common metadata. In a surrounding graph,
use identity maps on both outer carriers. At the selected edge use the related
box-body witness; at every other edge use reflexivity of recursive label
equivalence. The two outer incidence structures are identical, so their
preservation equations reduce to mapping a list by identity.
\end{proof}
\lean{Graph/RecursiveIsoPlug}{10}{19}{Linear.DeepIso.box}
\lean{Graph/RecursiveIsoPlug}{23}{42}{Linear.replaceBox_deepIso}

\begin{theorem}[Congruence at every addressed depth]
For a fixed host $G$ and typed site $s$, if $H\cong_r K$, then
\[
 G[s:=H]\cong_d G[s:=K].
\]
No context-uniqueness, freshness, or non-self-loop assumption is needed.
\end{theorem}
\begin{proof}
Induct on the site. A shallow site uses the generic shallow congruence with the
reflexive recursive-label relation. For a site inside a box, induction gives an
isomorphism between the two changed bodies. Lift it through the selected box
using box congruence. Repeating this step crosses every enclosing box in the
address.
\end{proof}
\lean{Graph/RecursiveIsoPlug}{46}{51}{Linear.plugAt_deepIso}

Two consequences connect this congruence with the preceding milestone. Exact
site freshness is invariant when the replacement is changed by recursive
isomorphism: the retained host predicate stays fixed, and the replacement hole
inventories coincide. Context validity of the two plugged results is also
equivalent, by context invariance applied to their constructed recursive
isomorphism. The latter statement does not assume the host is a context.
\lean{Graph/RecursiveIsoPlug}{54}{58}{Linear.siteFresh_deepIso_iff}
\lean{Graph/RecursiveIsoPlug}{61}{64}{Linear.plugAt_deepIso_unique_iff}

\subsection{Graphs and replacement modulo recursive renaming}

\begin{definition}[Recursive isomorphism classes]\label{term:recursive-quotient}
At fixed depth $d$ and interface $(I,O)$, declare $G\sim H$ when a recursive
isomorphism exists. Identity, inverse, and composition prove that this is an
equivalence relation, packaged as $\mathsf{deepSetoid}$. Its quotient is
$\mathsf{HypernetModuloIso}$, the type of graphs modulo independent carrier
renaming at all levels of their hierarchy. Write $[G]$ for the class of $G$.
\end{definition}
\lean{Graph/RecursiveIsoPlug}{67}{73}{Linear.deepSetoid}
\lean{Graph/RecursiveIsoPlug}{76}{77}{Linear.HypernetModuloIso}

\begin{definition}[Plugging a replacement class]\label{term:quotient-plugging}
For a fixed concrete host and site $s$, define
\[
 \mathsf{plugAtModuloIso}(s,[H])=[G[s:=H]].
\]
\end{definition}
\begin{proof}[Well-definedness]
If $H$ and $K$ represent the same source class, their recursive isomorphism
induces one between the plugged graphs by deep congruence. The results therefore
belong to the same target class. Quotient elimination defines the displayed
operation from this representative-independent function. Its value on a
concrete representative is the displayed equation by the computation rule of
quotient elimination.
\end{proof}
\lean{Graph/RecursiveIsoPlug}{81}{84}{Linear.plugAtModuloIso}
\lean{Graph/RecursiveIsoPlug}{87}{90}{Linear.plugAtModuloIso_mk}
\termdef{quotient-context-validity}
Context validity descends to classes as well: define
$\mathsf{moduloIsoUniqueHoles}([G])$ to mean that $G$ has unique holes. The
isomorphism-invariance theorem makes these propositions equivalent for any two
representatives; propositional extensionality supplies the equality required by
quotient elimination.
\lean{Graph/RecursiveIsoPlug}{93}{94}{Linear.moduloIsoUniqueHoles}

These constructions quotient the replacement argument for a fixed host and
address. Section~\ref{sec:site-transport} extends invariance to corresponding sites in
isomorphic hosts. Section~\ref{sec:composition} proves associativity and units.
Section~\ref{sec:deep-interchange} constructs the quotient of host/site pairs.

\subsection{A strict extension of literal-label isomorphism}

Use the earlier two-hole chain. The permutation $i\mapsto 1-i$ on its two-edge
carrier exchanges edge names. Applying it twice is identity because
$0\leq i<2$. Reindex the chain's edges by this bijection and keep its internal
vertex name unchanged; call the result $H'$. The existing reindexing isomorphism,
embedded into recursive isomorphism, witnesses $H\cong_0 H'$.
\lean{Graph/RecursiveIsoExamples}{9}{13}{Linear.Examples.swapTwo}
\lean{Graph/RecursiveIsoExamples}{15}{16}{Linear.Examples.renamedTwoHoles}
\lean{Graph/RecursiveIsoExamples}{18}{19}{Linear.Examples.twoHoles_deepIso}
The records $H$ and $H'$ are unequal. If they were equal, their edge-label
functions would agree at the edge named zero. Reading the hole name there gives
$0$ in $H$ and $1$ in $H'$, a contradiction.
\lean{Graph/RecursiveIsoExamples}{22}{29}{Linear.Examples.renamedTwoHoles_ne}

Box $H'$ with zero bound variables and one free output. Box congruence relates
this box to the earlier box of $H$.
\lean{Graph/RecursiveIsoExamples}{31}{32}{Linear.Examples.renamedBoxedTwoHoles}
\lean{Graph/RecursiveIsoExamples}{35}{36}{Linear.Examples.boxedTwoHoles_deepIso}
No literal label-preserving isomorphism relates these two boxed graphs. Such an
isomorphism's label equation at the sole source edge would make their box labels
equal. Equality of these labels forces equality of the body records, contrary
to the preceding result. This proves that recursive isomorphism identifies
strictly more graphs than the earlier literal-label notion.
\lean{Graph/RecursiveIsoExamples}{40}{46}{Linear.Examples.boxedTwoHoles_no_literalIso}
The renamed box remains a context by recursive invariance of unique holes. The
two different box records represent the same quotient class by quotient
soundness applied to their recursive isomorphism.
\lean{Graph/RecursiveIsoExamples}{48}{49}{Linear.Examples.renamedBoxedTwoHoles_unique}
\lean{Graph/RecursiveIsoExamples}{52}{55}{Linear.Examples.renamedBoxes_same_class}

For a nontrivial addressed replacement, use the earlier site selecting the first
hole inside a box. Substituting $H$ or $H'$ there gives recursively isomorphic
results by deep congruence. Finally, the primitive operations named
$\mathsf{value}$ and $\mathsf{otherValue}$ are not equivalent: at depth zero
recursive label equivalence is equality, and these are distinct constructors.
Thus carrier renaming does not erase semantic operation labels.
\lean{Graph/RecursiveIsoExamples}{58}{60}{Linear.Examples.deep_plug_renaming}
\lean{Graph/RecursiveIsoExamples}{63}{66}{Linear.Examples.distinct_operations_not_equivalent}
